% ============================================================================
%  D/17-ket-hop-tu — file chương
%  Ngày tạo: 2026-09-07 | Sửa gần nhất: 2026-09-07 | Ôn gần nhất: chưa
%  Dependency: A/01, A/02, D/16. Mức độ: cốt lõi.
% ============================================================================
\documentclass[../../english.tex]{subfiles}
\begin{document}
\chapter{Kết hợp từ: cái không có luật (Collocation)}
\ptrtarget{D/17}\ptrtarget{D/17-ket-hop-tu}
\dependency{\ptr{A/01} vì đây là hệ quả trực tiếp của ``nghĩa nằm ở cách dùng''; \ptr{D/16} cho chiều phạm vi áp dụng.}

\begin{tomtat}
Chương này nói về lớp ngôn ngữ mà người học thường bỏ qua lâu nhất và trả giá nhiều nhất: từ nào quen đi với từ nào. \emph{Make a decision} đúng, \emph{do a decision} sai --- không luật ngữ pháp nào giải thích được, chỉ có thói quen của cộng đồng nói. Nội dung gồm: kết hợp từ là gì và vì sao nó không phải ngữ pháp cũng không phải nghĩa; sáu kiểu kết hợp thường gặp; vì sao lỗi kết hợp lại gây phản ứng mạnh hơn lỗi ngữ pháp; các cụm đóng gói sẵn trong văn học thuật; và cách học chúng mà không phải thuộc lòng danh sách. Nếu chỉ nhớ một điều: đừng học từ đơn lẻ --- học \emph{cụm}, vì đơn vị mà người bản ngữ thật sự lưu trong đầu là cụm chứ không phải từ.
\end{tomtat}

\begin{nentang}
\kn{kết hợp từ}{hai hoặc nhiều từ hay xuất hiện cùng nhau hơn mức ngẫu nhiên: \emph{heavy rain}, \emph{make a decision}, \emph{strongly suggest}.}
\kn{cụm đóng gói}{chuỗi từ được lưu và dùng như một đơn vị: \emph{to the best of our knowledge}, \emph{in the context of}.}
\kn{độ chặt}{mức cố định của một tổ hợp, từ tự do (\emph{a big house}) tới cố định hoàn toàn (\emph{by and large}).}
\kn{lỗi chuyển di}{lỗi do dịch thẳng tổ hợp từ tiếng mẹ đẻ: \emph{open the light} từ ``mở đèn''.}
\kn{nguyên tắc chọn cụm}{giả thuyết rằng người nói lấy sẵn các cụm đã đóng gói thay vì ráp từng từ theo luật.}
\kn{học theo cụm}{ghi và ôn tập ở đơn vị cụm chứ không phải từ đơn.}
\end{nentang}

\begin{khoigoi}
\h{Vì sao \emph{do a decision} sai, dù không luật ngữ pháp nào cấm nó?}
\h{Vì sao lỗi kết hợp từ lại làm người bản ngữ khó chịu hơn lỗi chia động từ?}
\h{Nếu kết hợp từ không có luật, học nó thế nào cho khả thi --- chẳng lẽ thuộc lòng hàng vạn cặp?}
\h{Người bản ngữ nói trôi chảy trong thời gian thực bằng cách nào, nếu phải ráp từng từ theo luật?}
\h{Trong viết học thuật, những cụm nào đáng học trước nhất?}
\end{khoigoi}

Có một trải nghiệm quen thuộc với người học tiếng Anh ở trình độ khá: bạn viết một câu, kiểm lại ngữ pháp --- đúng hết, mọi từ đều tra từ điển --- và người bản ngữ vẫn nói ``nghe lạ''. Khi bạn hỏi sai chỗ nào, họ thường không giải thích được; họ chỉ nói ``chúng tôi không nói thế''.

Đó không phải sự né tránh. Họ thật sự không có luật để dẫn ra, vì cái bạn vi phạm không phải luật mà là \emph{thói quen}. Và thói quen này là một lớp ngôn ngữ riêng, tách khỏi cả ngữ pháp lẫn nghĩa.

\section{Không phải ngữ pháp, không phải nghĩa}

Câu hỏi mở đầu thứ nhất. Hãy xét \emph{do a decision}. Về ngữ pháp: \emph{do} là ngoại động từ, \emph{a decision} là danh từ đếm được có mạo từ --- hoàn toàn hợp lệ. Về nghĩa: \emph{do} nghĩa là làm, và ra quyết định là một việc bạn làm --- hoàn toàn hợp lý. Câu này sai theo một chiều thứ ba: cộng đồng nói tiếng Anh dùng \emph{make} với \emph{decision}, và chỉ vậy thôi.

Điều quan trọng cần thấy: sự phân công giữa \emph{make} và \emph{do} không suy ra được từ nghĩa của chúng. Bạn \emph{make} a decision, a mistake, an assumption, progress; bạn \emph{do} research, an experiment, a calculation. Nếu có logic, nó rất mờ và có quá nhiều ngoại lệ để dùng được. Cách duy nhất là biết.

\begin{tuvung}[Cùng nghĩa, khác chỗ dùng]
\tv{heavy}{Nặng}{Đi với \emph{rain}, \emph{traffic}, \emph{loss}, \emph{criticism}}{heavy rain, not strong rain}
\tv{strong}{Mạnh}{Đi với \emph{wind}, \emph{evidence}, \emph{correlation}, \emph{argument}}{strong evidence, not heavy evidence}
\tv{high}{Cao}{Đi với \emph{temperature}, \emph{cost}, \emph{accuracy}, \emph{risk}}{high accuracy, not tall accuracy}
\tv{deep}{Sâu}{Đi với \emph{understanding}, \emph{concern}, \emph{sleep}}{a deep understanding, not a high understanding}
\tv{make}{Làm, tạo}{\emph{a decision}, \emph{an assumption}, \emph{a mistake}, \emph{progress}}{make an assumption}
\tv{do}{Làm, thực hiện}{\emph{research}, \emph{an experiment}, \emph{a calculation}}{do an experiment}
\tv{draw}{Rút ra}{\emph{a conclusion}, \emph{a comparison}, \emph{an analogy}}{draw a conclusion, not take a conclusion}
\tv{pose}{Đặt ra}{\emph{a problem}, \emph{a question}, \emph{a challenge}}{This poses a challenge for \ldots}
\end{tuvung}

\section{Sáu kiểu thường gặp}

Biết các kiểu giúp bạn \emph{để ý} đúng chỗ khi đọc --- đó là mục đích của mục này, không phải để thuộc tên gọi.

\emph{Một --- động từ \(+\) danh từ.} \emph{Make an assumption}, \emph{draw a conclusion}, \emph{conduct an experiment}, \emph{address a limitation}. Kiểu này gây lỗi nhiều nhất cho người học và cũng dễ sửa nhất.

\emph{Hai --- tính từ \(+\) danh từ.} \emph{Strong evidence}, \emph{high accuracy}, \emph{significant improvement}, \emph{substantial gap}.

\emph{Ba --- trạng từ \(+\) tính từ hoặc động từ.} \emph{Strongly suggest}, \emph{widely used}, \emph{closely related}, \emph{highly sensitive}. Kiểu này quan trọng đặc biệt cho \ptr{C/14}, vì trạng từ là công cụ chính điều chỉnh mức cam kết.

\emph{Bốn --- danh từ \(+\) danh từ.} \emph{Error rate}, \emph{sensitivity analysis}, \emph{ground truth}. Trong tiếng Anh kỹ thuật, các chuỗi danh từ dài là chuẩn --- và cũng là chỗ dễ viết ra thứ không ai dùng.

\emph{Năm --- động từ \(+\) giới từ.} \emph{Depend on}, \emph{result in}, \emph{account for}, \emph{consist of}. Giới từ ở đây gần như không mang nghĩa; nó chỉ là phần cố định của cụm, và không dịch được.

\emph{Sáu --- cụm đóng gói dài.} \emph{To the best of our knowledge}, \emph{it is worth noting that}, \emph{in the context of}, \emph{with respect to}. Đây là kiểu có tỉ lệ lợi ích trên công sức cao nhất trong viết học thuật --- xem mục 5.

\section{Vì sao lỗi kết hợp gây phản ứng mạnh}

Câu hỏi mở đầu thứ hai có một câu trả lời liên quan tới cách người nghe xử lý ngôn ngữ, và nó cũng trả lời luôn câu hỏi mở đầu thứ tư.

Giả thuyết được nhiều nhà nghiên cứu ngữ liệu ủng hộ --- Sinclair phát biểu rõ nhất\cite{sinclair1991} --- là: người nói không ráp từng từ theo luật ngữ pháp trong thời gian thực. Phần lớn thời gian họ lấy ra các \emph{cụm đã đóng gói sẵn} và chỉ điều chỉnh nhẹ. Ngữ pháp được dùng để nối các cụm đó lại và để xử lý các trường hợp mới.

Điều này giải thích được một chuyện khó hiểu: làm sao người nói tạo được câu trôi chảy nhanh đến vậy, khi việc lập kế hoạch từng từ theo luật là quá chậm. Câu trả lời: phần lớn công việc đã được làm sẵn và lưu thành khối.

Từ đó ra hệ quả cho người nghe. Khi bạn dùng đúng cụm, người nghe nhận ra khối quen thuộc và xử lý gần như không tốn công. Khi bạn dùng \emph{do a decision}, khối đó không khớp với cái nào trong đầu họ, và họ phải chuyển sang xử lý từng từ --- chậm hơn nhiều. Cảm giác chủ quan của họ là ``vấp''.

Còn lỗi chia động từ thì sao? Nó thường \emph{không} phá khối: \emph{he make a decision} vẫn khớp với khối \emph{make a decision}, chỉ sai một hình thái mà người nghe tự sửa. Đó là lý do một lỗi mà sách giáo khoa coi là nghiêm trọng lại ít gây vấp hơn một lỗi mà sách giáo khoa không nhắc tới.

\begin{figure}[H]\centering
\begin{tikzpicture}[
  blk/.style={draw=steel,fill=bluebg,rounded corners=2pt,inner sep=4pt,font=\footnotesize,align=center,minimum height=0.62cm},
  bad/.style={blk,draw=reddark,fill=redbg},
  lb/.style={font=\scriptsize\itshape,color=tiercol,align=left,text width=3.5cm}]
\node[blk,text width=3.2cm] (a1) {To the best of our knowledge};
\node[blk,right=0.12cm of a1,text width=1.5cm] (a2) {, no prior};
\node[blk,right=0.12cm of a2,text width=2.4cm] (a3) {work has addressed};
\node[lb,right=0.28cm of a3] {3 khối \(\rightarrow\) xử lý nhanh, không vấp};
\node[bad,below=0.9cm of a1.south west,anchor=north west,text width=1.0cm] (b1) {To};
\node[bad,right=0.10cm of b1,text width=0.8cm] (b2) {the};
\node[bad,right=0.10cm of b2,text width=0.9cm] (b3) {best};
\node[bad,right=0.10cm of b3,text width=0.6cm] (b4) {in};
\node[bad,right=0.10cm of b4,text width=0.7cm] (b5) {our};
\node[bad,right=0.10cm of b5,text width=1.6cm] (b6) {knowledges};
\node[lb,right=0.28cm of b6,text width=3.5cm] {khối vỡ \(\rightarrow\) phải ráp từng từ};
\end{tikzpicture}
\caption{Vì sao lỗi kết hợp gây vấp. Một giới từ sai (\emph{in} thay vì \emph{of}) phá vỡ khối quen thuộc, buộc người đọc chuyển sang xử lý từng từ --- tốn công hơn nhiều so với tự sửa một lỗi chia động từ.}
\label{fig:khoi-cum}
\end{figure}

\section{Lỗi chuyển di từ tiếng mẹ đẻ}

Nguồn lỗi kết hợp lớn nhất không phải thiếu vốn từ mà là dịch thẳng tổ hợp từ tiếng mẹ đẻ. Bạn biết cả hai từ, và bạn ghép chúng theo cách tiếng Việt ghép.

Ví dụ kinh điển ở nhiều ngôn ngữ: ``mở đèn'' thành \emph{open the light} thay vì \emph{turn on the light}. Nhưng các trường hợp nguy hiểm hơn nằm ở văn học thuật, nơi lỗi tinh vi hơn và bạn không tự phát hiện được: ``đưa ra kết luận'' thành \emph{give a conclusion} thay vì \emph{draw a conclusion}; ``giải quyết vấn đề'' thành \emph{solve the limitation} thay vì \emph{address the limitation}.

Cách phòng thực dụng: khi viết một tổ hợp mà bạn dịch từ tiếng Việt, hãy đánh dấu nó và kiểm bằng phép thử ba câu ở \ptr{D/16} --- tìm tổ hợp đó trong tài liệu ngành. Nếu không có, gần như chắc chắn nó là lỗi chuyển di.

\begin{cambay}
Đừng tin cảm giác ``nghe được''. Cảm giác của bạn về tiếng Anh được hình thành phần lớn từ tiếng Việt, nên một tổ hợp dịch thẳng sẽ nghe rất tự nhiên \emph{với bạn} và rất lạ với người đọc. Đây là lý do kết hợp từ là chỗ mà việc tra cứu bằng dữ liệu thật không thay thế được bằng trực giác --- kể cả khi trình độ của bạn đã khá.
\end{cambay}

\section{Cụm nào đáng học trước}

Câu hỏi mở đầu thứ năm. Trong viết học thuật, các cụm đóng gói kiểu sáu cho hiệu quả nhanh nhất, vì chúng lặp lại rất nhiều và mỗi cụm dùng được trong hàng chục bài.

\begin{mauau}[Cụm học thuật đáng thuộc]
\mc{Nêu giới hạn hiểu biết}{To the best of our knowledge, \ldots}{Dùng khi tuyên bố cái mới; an toàn hơn \emph{no one has}}
\mc{Hướng chú ý}{It is worth noting that \ldots}{Đánh dấu một chi tiết dễ bị bỏ qua}
\mc{Nêu điều kiện}{provided that \ldots\ / as long as \ldots}{Gọn hơn một mệnh đề \emph{if} dài}
\mc{Đặt trong bối cảnh}{in the context of \ldots}{Rất hay dùng; thay cho \emph{about} khi cần trang trọng}
\mc{Đối lập}{in contrast to \ldots\ / unlike \ldots}{\emph{In contrast to} cần một danh từ theo sau}
\mc{Giải thích lý do}{This is because \ldots\ / owing to \ldots}{\emph{Owing to} \(+\) danh từ, không \(+\) mệnh đề}
\mc{Chuyển sang phần mới}{We now turn to \ldots}{Báo đứt mạch có chủ ý (\ptr{C/12})}
\mc{Thừa nhận rồi phản bác}{While X holds in general, \ldots}{Nước đi 2 của CARS (\ptr{C/13})}
\mc{Nêu hệ quả}{It follows that \ldots}{Mạnh; chỉ dùng khi thật sự suy ra được}
\mc{Giới hạn phạm vi}{within the scope of this work, \ldots}{Công cụ hai của \ptr{C/14}}
\end{mauau}

Cách học ba mươi cụm như thế này hiệu quả hơn nhiều so với học ba trăm từ đơn, vì chúng dùng được ngay, chúng đúng chắc chắn, và chúng giải phóng sự chú ý của bạn cho nội dung.

\section{Học kết hợp từ mà không thuộc lòng}

Câu hỏi mở đầu thứ ba là câu hỏi thực tế nhất: nếu không có luật, chẳng lẽ phải thuộc hàng vạn cặp? Không, vì bạn không cần hàng vạn --- bạn cần các cặp trong \emph{ngành bạn}, và số đó nhỏ hơn nhiều. Bốn nguyên tắc.

\emph{Một --- ghi ở đơn vị cụm, không phải từ.} Khi gặp từ mới, đừng ghi \emph{assumption --- giả định}. Ghi \emph{make an assumption}, \emph{a strong assumption}, \emph{relax an assumption}. Đây là thay đổi nhỏ trong thói quen ghi chép nhưng thay đổi lớn trong kết quả, và nó khớp với trường ``cụm đi kèm'' trong bảng thuật ngữ ở \ptr{B/10}.

\emph{Hai --- đọc để ý cụm, không chỉ để ý nghĩa.} Khi đọc bài báo, thỉnh thoảng dừng lại ở một câu và hỏi: tác giả dùng động từ nào với danh từ này? Vài phút mỗi ngày, tích luỹ nhanh hơn bạn tưởng.

\emph{Ba --- dùng ngữ liệu ngành để kiểm, không để học.} Ngữ liệu tốt cho việc \emph{xác nhận} một tổ hợp bạn đang phân vân; nó kém cho việc học chủ động vì quá nhiều dữ liệu.

\emph{Bốn --- tái sử dụng có ý thức.} Khi bạn thấy một cụm hay trong bài báo, hãy chép nó vào một tệp riêng --- một ``ngân hàng cụm'' của bạn --- và cố ý dùng nó trong bài kế tiếp. Đây là cách chuyển từ nhận ra sang dùng được, và nó là dạng luyện có chủ đích cho ngôn ngữ.

\begin{cannho}
Ranh giới đạo văn: tái sử dụng \emph{cụm chức năng} (\emph{it is worth noting that}, \emph{in contrast to}) là hoàn toàn bình thường và được mọi cộng đồng chấp nhận --- chúng là công cụ chung, không thuộc về ai. Tái sử dụng \emph{câu mang nội dung} của người khác thì không, kể cả khi bạn đổi vài từ. Phép thử: nếu cụm đó có thể xuất hiện trong bất kỳ bài nào của ngành, nó là công cụ chung; nếu nó chứa ý riêng của tác giả, nó không phải.
\end{cannho}

\section{Gốc rễ: từ ``bạn nhận ra một từ qua bạn bè của nó''}

Ý trung tâm của chương này được Firth phát biểu năm 1957 trong một câu thường được trích\cite{firth1957}: bạn sẽ nhận ra một từ qua những từ nó đi cùng. Câu này nghe như một nhận xét văn vẻ, nhưng nó là một tuyên bố phương pháp luận: nếu nghĩa nằm ở cách dùng, thì muốn nghiên cứu nghĩa phải \emph{đếm} cách dùng, và muốn đếm thì cần kho văn bản lớn.

Vào thời điểm đó việc đó gần như bất khả thi bằng tay. Điều thay đổi mọi thứ là máy tính. Từ những năm 1960 và mạnh dần về sau, các kho ngữ liệu điện tử cho phép đếm chính xác từ nào đi với từ nào và tần suất bao nhiêu --- và kết quả gây bất ngờ cho chính giới ngôn ngữ học: mức độ mà ngôn ngữ thật bị chi phối bởi các tổ hợp quen thuộc lớn hơn nhiều so với hình dung của các mô tả ngữ pháp truyền thống.

Sinclair rút ra từ đó một phát biểu có ảnh hưởng lớn\cite{sinclair1991}: ngôn ngữ được tạo ra theo hai nguyên tắc song song. Nguyên tắc thứ nhất là ráp từ theo luật ngữ pháp --- linh hoạt nhưng chậm. Nguyên tắc thứ hai là lấy sẵn các cụm đã đóng gói --- nhanh và chiếm phần lớn. Các mô tả ngữ pháp truyền thống chỉ mô tả nguyên tắc thứ nhất, nên chúng bỏ sót đúng phần chiếm đa số.

Điều này có hệ quả trực tiếp cho cách bạn học. Nếu ngôn ngữ chủ yếu là các khối đóng gói, thì việc học ``ngữ pháp cộng từ vựng'' --- luật ở một bên, danh sách từ ở bên kia --- bỏ sót lớp ở giữa, đúng lớp mà người bản ngữ thật sự dùng. Các từ điển kết hợp từ và các từ điển cho người học dựa trên ngữ liệu, xuất hiện từ cuối những năm 1980\cite{cobuild1987,ldoce}, chính là phản ứng với chỗ trống đó. Cách dạy đặt cụm vào trung tâm thay vì đặt ngữ pháp vào trung tâm cũng ra đời trong giai đoạn này\cite{lewis1993}, và về sau được đặt trên nền dữ liệu về vai trò của các đơn vị đóng gói sẵn trong ngôn ngữ\cite{wray2002} cùng ý tưởng rằng mỗi từ mang theo một bộ mồi dùng riêng của nó\cite{hoey2005}.

Một điểm nữa đáng nhắc vì nó nối chương này với công nghệ hiện tại: các mô hình ngôn ngữ hiện đại học chính xác cái mà Firth mô tả --- dự đoán một từ từ các từ xung quanh nó. Việc chúng nắm rất tốt kết hợp từ không phải tình cờ; đó là hệ quả trực tiếp của giả thuyết nền.

\begin{gocre}
\gr{1957 --- Firth\cite{firth1957}}{Nghiên cứu nghĩa bằng cách nào?}{``Nhận ra một từ qua bạn bè của nó'' \(\Rightarrow\) muốn biết nghĩa thì đếm cách dùng.}
\gr{Từ 1960s --- ngữ liệu điện tử}{Đếm bằng tay không nổi.}{Máy tính cho phép đo tần suất tổ hợp trên kho lớn.}
\gr{Phát hiện gây bất ngờ}{Ngôn ngữ thật giống mô tả ngữ pháp tới đâu?}{Tỉ lệ tổ hợp quen thuộc cao hơn nhiều so với hình dung trước đó.}
\gr{Sinclair --- hai nguyên tắc\cite{sinclair1991}}{Vì sao người nói nhanh đến vậy?}{Lấy sẵn cụm đóng gói song song với ráp theo luật; ngữ pháp truyền thống chỉ mô tả một nửa.}
\gr{Cuối 1980s --- từ điển dựa ngữ liệu\cite{cobuild1987}}{Chỗ trống giữa ngữ pháp và danh sách từ?}{Từ điển cho người học ghi cụm đi kèm và nhãn văn vực.}
\gr{Mô hình ngôn ngữ hiện đại}{Máy học ngôn ngữ thế nào?}{Dự đoán từ từ ngữ cảnh --- đúng giả thuyết của Firth, nên chúng nắm kết hợp từ rất tốt.}
\end{gocre}

\section{Liên hệ ngang}

\begin{lienhe}
\lh{Học máy}{Mô hình ngôn ngữ và nhúng từ}{Hiện thực hoá trực tiếp giả thuyết Firth: nghĩa suy từ ngữ cảnh. Giải thích vì sao chúng giỏi kết hợp từ và có ích cho việc kiểm ở mục 6.}
\lh{Chơi cờ}{Trí nhớ theo khối của kỳ thủ}{Kỳ thủ mạnh nhớ thế cờ theo khối có ý nghĩa chứ không theo từng quân --- cùng cơ chế với việc lưu cụm ở mục 3.}
\lh{Âm nhạc}{Ngón đàn và mô-típ}{Nhạc công không nghĩ từng nốt mà lấy sẵn các mô-típ đã luyện; ứng biến là ghép và biến đổi khối.}
\lh{Lập trình}{Thành ngữ của một ngôn ngữ}{Mã ``đúng nhưng lạ'' là hiện tượng y hệt: hợp cú pháp mà không theo cách cộng đồng viết. Người rà soát phản ứng giống người bản ngữ.}
\lh{Tiếng Việt}{Từ ghép và cụm cố định}{``Đưa ra quyết định'' chứ không ``làm quyết định''; tiếng Việt có cùng lớp thói quen, và chính nó gây lỗi chuyển di ở mục 4.}
\lh{Tâm lý học}{Bộ nhớ làm việc và tổ chức theo khối}{Nền tảng chung: gộp thông tin thành khối lớn hơn để vượt giới hạn bộ nhớ ngắn hạn --- cùng nguyên lý với \ptr{C/11} mục 3.}
\end{lienhe}

\begin{traloi}
\tl{Vì nó không vi phạm luật nào --- nó vi phạm \emph{thói quen}. Về ngữ pháp, \emph{do} là ngoại động từ và \emph{a decision} là danh từ đếm được có mạo từ: hợp lệ. Về nghĩa, ra quyết định đúng là việc bạn làm: hợp lý. Nhưng cộng đồng nói tiếng Anh dùng \emph{make} với \emph{decision}, và sự phân công giữa \emph{make} và \emph{do} không suy ra được từ nghĩa của chúng. Đây là một lớp ngôn ngữ riêng, tách khỏi cả ngữ pháp lẫn nghĩa.}
\tl{Vì nó phá vỡ khối. Người nghe không xử lý từng từ mà nhận ra các cụm đã đóng gói sẵn; khi cụm khớp, việc xử lý gần như không tốn công. \emph{Do a decision} không khớp khối nào nên họ phải chuyển sang ráp từng từ --- chậm hơn nhiều, và cảm giác chủ quan là ``vấp''. Ngược lại, \emph{he make a decision} vẫn khớp khối \emph{make a decision}, chỉ sai một hình thái mà người nghe tự sửa --- nên lỗi mà sách giáo khoa coi là nghiêm trọng lại ít gây vấp hơn.}
\tl{Bạn không cần hàng vạn cặp --- chỉ cần các cặp trong ngành bạn, và số đó nhỏ hơn nhiều. Bốn nguyên tắc: ghi ở đơn vị cụm chứ không phải từ (\emph{make an assumption} thay vì \emph{assumption --- giả định}); khi đọc thì thỉnh thoảng dừng lại hỏi tác giả dùng động từ nào với danh từ này; dùng ngữ liệu để \emph{kiểm} một tổ hợp đang phân vân chứ không để học; và giữ một ngân hàng cụm riêng rồi cố ý dùng chúng trong bài kế tiếp.}
\tl{Bằng cách không ráp từng từ. Giả thuyết được ngữ liệu ủng hộ là ngôn ngữ được tạo theo hai nguyên tắc song song: ráp theo luật ngữ pháp --- linh hoạt nhưng chậm --- và lấy sẵn các cụm đã đóng gói --- nhanh và chiếm phần lớn. Ngữ pháp được dùng để nối các cụm và xử lý trường hợp mới. Điều này giải thích tốc độ nói mà việc lập kế hoạch từng từ không giải thích được, và nó cũng nói rằng mô tả ``ngữ pháp cộng từ vựng'' bỏ sót lớp chiếm đa số.}
\tl{Các cụm đóng gói dài trong văn học thuật, vì chúng lặp lại rất nhiều và mỗi cụm dùng được trong hàng chục bài: \emph{to the best of our knowledge}, \emph{it is worth noting that}, \emph{in the context of}, \emph{in contrast to}, \emph{it follows that}, \emph{within the scope of this work}. Học ba mươi cụm như vậy hiệu quả hơn ba trăm từ đơn --- chúng dùng được ngay, đúng chắc chắn, và giải phóng chú ý cho nội dung. Tái sử dụng cụm chức năng là bình thường; tái sử dụng câu mang nội dung thì không.}
\end{traloi}

\begin{dongian}
Có một lớp tiếng Anh mà sách giáo khoa gần như không dạy, nhưng nó lại là thứ khiến câu của bạn nghe ``lạ'' dù không sai chữ nào.

Đó là chuyện từ nào quen đi với từ nào. Bạn nói \emph{make a decision} chứ không nói \emph{do a decision}. \emph{Heavy rain} chứ không \emph{strong rain}. \emph{Strong evidence} chứ không \emph{heavy evidence}. Không có luật nào giải thích được --- người bản ngữ cũng không giải thích được, họ chỉ nói ``chúng tôi không nói thế''.

Vì sao lỗi kiểu này lại khó chịu hơn cả lỗi chia động từ? Vì người đọc không xử lý từng từ một. Họ nhận ra cả \emph{cụm} quen thuộc và lướt qua rất nhanh. Khi bạn viết \emph{do a decision}, cụm đó không có trong đầu họ, nên họ phải dừng lại và ghép từng từ. Còn \emph{he make a decision} thì cụm vẫn nguyên, chỉ thiếu một chữ ``s'' mà họ tự thêm vào.

Nguồn lỗi lớn nhất là dịch thẳng từ tiếng Việt. ``Đưa ra kết luận'' thành \emph{give a conclusion} --- trong khi tiếng Anh nói \emph{draw a conclusion}. Nguy hiểm ở chỗ những câu này nghe rất tự nhiên \emph{với bạn}, vì cảm giác của bạn được hình thành từ tiếng Việt.

Vậy học thế nào? Không phải thuộc lòng hàng vạn cặp. Chỉ cần một thay đổi nhỏ trong cách ghi chép: \emph{đừng ghi từ, hãy ghi cụm}. Thay vì ghi ``assumption = giả định'', hãy ghi \emph{make an assumption}, \emph{a strong assumption}, \emph{relax an assumption}. Cùng công sức, nhưng cái bạn ghi là thứ dùng được ngay.

Và nếu muốn tiến bộ nhanh trong viết học thuật, hãy học khoảng ba mươi \emph{cụm dài} sẵn có: \emph{to the best of our knowledge}, \emph{it is worth noting that}, \emph{in contrast to}, \emph{it follows that}. Chúng lặp lại trong mọi bài báo, chúng luôn đúng, và dùng chúng giải phóng đầu óc bạn cho phần quan trọng hơn --- nội dung.
\motcau{Học cụm, đừng học từ.}
\chuadongian{Biết một cụm là công cụ chung hay là ý riêng của tác giả đòi bạn quen với ngành --- đó là ranh giới giữa tái sử dụng và đạo văn.}
\end{dongian}

\begin{onbai}
\ar[3]{Kết hợp từ là gì}{Từ nào quen đi với từ nào; không phải ngữ pháp, không phải nghĩa, mà là thói quen cộng đồng.}
\ar[3]{Vì sao lỗi kết hợp gây vấp}{Nó phá vỡ khối quen thuộc, buộc người đọc chuyển sang xử lý từng từ.}
\ar[3]{Vì sao lỗi chia động từ ít gây vấp hơn}{Khối vẫn nguyên; người đọc tự sửa hình thái.}
\ar[3]{Nguyên tắc học chính}{Ghi ở đơn vị cụm, không phải từ đơn.}
\ar[3]{Sáu kiểu kết hợp}{Động từ\(+\)danh từ; tính từ\(+\)danh từ; trạng từ\(+\)tính từ/động từ; danh từ\(+\)danh từ; động từ\(+\)giới từ; cụm đóng gói dài.}
\ar[2]{Ví dụ phân công không có logic}{\emph{make} a decision/an assumption; \emph{do} research/an experiment.}
\ar[2]{Hai nguyên tắc tạo ngôn ngữ}{Ráp theo luật (linh hoạt, chậm) và lấy sẵn cụm đóng gói (nhanh, chiếm phần lớn).}
\ar[2]{Lỗi chuyển di}{Dịch thẳng tổ hợp tiếng Việt: \emph{give a conclusion} thay vì \emph{draw a conclusion}.}
\ar[2]{Vì sao không tin được cảm giác}{Cảm giác ``nghe được'' của bạn hình thành từ tiếng Việt, nên tổ hợp dịch thẳng nghe tự nhiên với bạn.}
\ar[2]{Cụm đáng học trước}{Các cụm đóng gói học thuật --- lặp nhiều, luôn đúng, dùng được ngay.}
\ar[2]{Ranh giới đạo văn}{Cụm chức năng dùng chung được; câu mang nội dung riêng thì không.}
\ar[1]{Nguồn gốc ý tưởng}{Firth 1957: nhận ra một từ qua những từ nó đi cùng \(\Rightarrow\) muốn biết nghĩa phải đếm cách dùng.}
\end{onbai}

\begin{hoidap}
\qa{Có từ điển riêng cho kết hợp từ không?}{Có, và chúng là loại từ điển bị dùng ít nhất so với mức hữu ích. Chúng liệt kê, với mỗi danh từ, những động từ và tính từ hay đi cùng --- đúng thông tin mà từ điển thường không cho. Ngoài ra, các từ điển tốt cho người học cũng có mục cụm đi kèm ở cuối mỗi từ. Nếu bạn chỉ thêm một công cụ vào thói quen tra cứu, hãy thêm cái này\cite{ldoce,manchesterphrasebank}.}
\qa{Tôi nên học bao nhiêu cụm học thuật?}{Ba mươi tới năm mươi cụm chức năng là đủ cho phần lớn nhu cầu viết bài, vì chúng lặp lại rất nhiều. Quan trọng hơn số lượng là cách chọn: lấy chúng từ chính các bài trong ngành bạn thay vì từ danh sách chung, vì mỗi cộng đồng có bộ cụm ưa dùng hơi khác nhau. Cách làm: khi đọc, đánh dấu mọi cụm chức năng bạn thấy, và sau mười bài bạn sẽ thấy khoảng ba mươi cái lặp đi lặp lại --- đó là danh sách của bạn.}
\qa{Học cụm nhiều có làm bài của tôi nghe rập khuôn không?}{Không, vì các cụm chức năng nằm ở phần \emph{khung} của câu chứ không ở phần nội dung. Cái làm bài bạn khác biệt là ý, số liệu, và cách lập luận --- không phải việc bạn viết \emph{it is worth noting that} hay \emph{note that}. Thực tế thì ngược lại: dùng cụm chuẩn giải phóng chú ý của bạn cho nội dung, nên bài thường mạnh hơn chứ không nhạt đi.}
\qa{Mô hình ngôn ngữ có tin được cho việc kiểm kết hợp từ không?}{Đây là một trong những việc chúng làm đáng tin nhất, vì chúng được huấn luyện chính xác trên hiện tượng này. Cách hỏi hiệu quả: đưa cả câu và hỏi ``tổ hợp này có tự nhiên trong văn viết kỹ thuật không, nếu không thì người ta thường nói thế nào'' --- hỏi so sánh cho kết quả tốt hơn hỏi đúng/sai. Vẫn nên kiểm chéo bằng tài liệu ngành cho các tổ hợp mang thuật ngữ, vì đó là chỗ tiếng Anh chung và tiếng Anh ngành khác nhau nhất.}
\qa{Đây có phải là thứ chỉ cần cho người viết, hay cả người đọc?}{Cả hai, và tác dụng khi đọc ít được nói tới. Nhận ra cụm khiến bạn đọc nhanh hơn hẳn, vì bạn xử lý khối thay vì từng từ --- đúng cơ chế ở mục 3. Ngoài ra, một số cụm mang tín hiệu về thái độ và mức cam kết của tác giả (\ptr{B/08}): \emph{it should be noted that} thường báo trước một điều chỉnh hoặc một giới hạn quan trọng. Học cụm vì thế cải thiện cả tốc độ lẫn độ chính xác khi đọc.}
\qa{Tôi hay quên cụm đã ghi. Làm sao nhớ được?}{Vấn đề thường không phải trí nhớ mà là bạn ghi mà không dùng. Cách chữa hiệu quả nhất là nguyên tắc bốn ở mục 6: mỗi lần viết một đoạn, cố ý dùng ba cụm từ ngân hàng của bạn. Một cụm được dùng hai ba lần trong bài thật sẽ ở lại; một cụm chỉ nằm trong sổ thì không. Về hệ thống ôn tập lặp lại, \ptr{E/19} và \ptr{E/20} nói kỹ hơn, nhưng nguyên tắc chung là ưu tiên dùng hơn ôn.}
\qa{Trong tiếng Việt tôi có kết hợp từ không, hay chỉ tiếng Anh mới có?}{Mọi ngôn ngữ đều có, và tiếng Việt cũng vậy: ``đưa ra quyết định'' chứ không ``làm quyết định''; ``rút ra kết luận'' chứ không ``kéo ra kết luận''. Bạn không để ý vì bạn đã thuộc chúng từ nhỏ. Nhận ra điều này có ích thực tế: nó cho bạn thấy lớp này tồn tại và có thể học được, và nó giải thích chính xác vì sao dịch thẳng lại hỏng --- hai ngôn ngữ có hai bộ thói quen không trùng nhau.}
\end{hoidap}

\begin{baitap}
\bt[1]{Chọn 10 danh từ hay dùng trong ngành bạn và tìm động từ đi với mỗi từ}{Tài liệu ngành}{Ghi cả cụm, không ghi từ đơn.}
\bt[1]{Đánh dấu mọi cụm chức năng trong một bài báo}{Bài báo ngành}{Đếm xem có bao nhiêu cụm khác nhau; thường ít hơn bạn tưởng.}
\bt[2]{Lập ngân hàng 30 cụm học thuật từ 10 bài trong ngành bạn}{Tài liệu ngành}{Chỉ giữ cụm xuất hiện ở ít nhất hai bài khác nhau.}
\bt[2]{Tìm 10 tổ hợp trong bài của bạn có thể là lỗi chuyển di, và kiểm bằng phép thử ba câu}{Bài viết của bạn}{Đánh dấu trước những chỗ bạn nhớ là đã dịch từ tiếng Việt.}
\bt[2]{Viết một đoạn 200 từ dùng ít nhất 5 cụm từ ngân hàng của bạn}{Bài viết của bạn}{Cố ý dùng, không chờ chúng tự đến.}
\bt[3]{Với 5 tính từ (\emph{strong}, \emph{high}, \emph{heavy}, \emph{deep}, \emph{significant}), lập danh sách danh từ đi với mỗi từ trong ngành bạn}{Tài liệu ngành}{Bạn sẽ thấy sự phân công rất rõ và không có logic.}
\bt[3]{Trong bốn tuần, mỗi tuần thêm 10 cụm vào ngân hàng và dùng 5 trong số đó}{Bài viết của bạn}{Đo bằng số cụm \emph{đã dùng}, không phải số cụm đã ghi.}
\end{baitap}

\section*{Kết chương và đọc tiếp}

Chương này cho một lớp ngôn ngữ mà cách học ``ngữ pháp cộng từ vựng'' bỏ sót, một giải thích vì sao lỗi ở lớp này gây phản ứng mạnh, và một thay đổi nhỏ trong thói quen ghi chép --- ghi cụm thay vì ghi từ --- có tác dụng lớn hơn hầu hết các bài luyện.

Hai chương vừa rồi đều nói về những gì quyết định nghĩa \emph{bên ngoài} định nghĩa từ điển: sắc thái ở mức từ, thói quen ở mức tổ hợp. Còn một mức nữa, và nó lớn hơn cả hai. Khi một ngành nói về đối tượng của mình, nó luôn nói qua một hệ hình ảnh mượn từ chỗ khác --- bộ nhớ như một cái kho, mạng thần kinh như một bộ não, thị trường như một sinh vật. Những hình ảnh đó không phải trang trí; chúng quyết định câu hỏi nào được coi là hợp lý và câu hỏi nào không ai nghĩ tới. Chương \ptr{D/18} chỉ ra chúng, kể cả trong chính ngành của bạn.
\ifSubfilesClassLoaded{\printbibliography[title={Nguồn}]}{}
\end{document}
